% ============================================================
% 6yr_data_cleaning_summary.tex
% Methods note: cleaning of EPA Six-Year Review contaminant
% concentration data for the inorganic-chemical analysis.
% Mirrors Ravalli et al. (2022, Lancet Planet. Health) appendix,
% applied to the full national panel before the downstream-CWS cut.
% Requires: booktabs. Include with % ============================================================
% 6yr_data_cleaning_summary.tex
% Methods note: cleaning of EPA Six-Year Review contaminant
% concentration data for the inorganic-chemical analysis.
% Mirrors Ravalli et al. (2022, Lancet Planet. Health) appendix,
% applied to the full national panel before the downstream-CWS cut.
% Requires: booktabs. Include with % ============================================================
% 6yr_data_cleaning_summary.tex
% Methods note: cleaning of EPA Six-Year Review contaminant
% concentration data for the inorganic-chemical analysis.
% Mirrors Ravalli et al. (2022, Lancet Planet. Health) appendix,
% applied to the full national panel before the downstream-CWS cut.
% Requires: booktabs. Include with % ============================================================
% 6yr_data_cleaning_summary.tex
% Methods note: cleaning of EPA Six-Year Review contaminant
% concentration data for the inorganic-chemical analysis.
% Mirrors Ravalli et al. (2022, Lancet Planet. Health) appendix,
% applied to the full national panel before the downstream-CWS cut.
% Requires: booktabs. Include with \input{output/6yr_data_cleaning_summary}.
% ============================================================

\section{Cleaning of EPA Six-Year Review concentration data}
\label{sec:6yr_cleaning}

Measured contaminant concentrations are drawn from the U.S.\ EPA Six-Year
Review compliance-monitoring database, pooling the Second Six-Year Review
(SYR2, sample years 1998--2005) and Third Six-Year Review (SYR3, 2006--2011).
We clean these records following the procedure of Ravalli et al.\
(2022),\footnote{Ravalli F, Yu Y, Bostick BC, et al.\ Sociodemographic
inequalities in uranium and other metals in community water systems across the
USA, 2006--11: a cross-sectional study. \emph{Lancet Planet Health} 2022;
\textbf{6}: e320--30 (Supplementary appendix~2).} with one deliberate
departure described in Section~\ref{subsec:mdl}. All cleaning is performed on
the \emph{full national} sample of community water systems---before the analysis
sample is restricted to systems located downstream of a coal mine---so that
detection-rate screening, limit-of-detection imputation, and outlier removal
exploit the entire monitoring record rather than the small downstream subsample.

\subsection{Cleaning steps}

\begin{enumerate}

  \item \textbf{Community water systems only.} We retain records from community
    water systems (CWS) and drop transient and non-transient non-community
    systems, consistent with our focus on systems serving residential
    populations year-round. This removes roughly 23\% of raw monitoring records.

  \item \textbf{Unit normalization.} All mass-concentration readings are
    converted to mg/L and all radionuclide readings to pCi/L, so that records
    reported in $\mu$g/L, ng/L, or other units are placed on a common scale
    before averaging.

  \item \textbf{Limit-of-detection substitution.} Non-detect samples---which the
    database reports with a missing concentration---are assigned a value equal
    to the reported limit of detection (LOD) divided by $\sqrt{2}$, the standard
    convention used by the CDC and other federal agencies when summarizing
    environmental concentration data. A record's own reported LOD is used when it
    is reported in mg/L or $\mu$g/L and does not exceed 5~$\mu$g/L (the reliable
    range); otherwise the LOD is taken from a per-chemical method detection limit
    (MDL), as described in Section~\ref{subsec:mdl}. Imputing non-detects, rather
    than discarding them, is the main reason the analysis sample roughly doubles
    in size relative to a detections-only sample.

  \item \textbf{Gross-outlier removal.} Records reporting a concentration greater
    than 100 times the contaminant's maximum contaminant level (MCL) are dropped
    as likely data-entry errors.

  \item \textbf{Detection-rate screen.} For each contaminant we compute the
    national share of records above the LOD and retain only contaminants
    detected in more than 10\% of records. Among the inorganic chemicals this
    keeps arsenic, nitrate, barium, chromium, and selenium, and excludes cadmium
    (1.8\% detected), mercury (1.6\%), and thallium, whose concentrations are too
    rarely above detection to support a credible mean-concentration regression.
    The threshold is a tunable parameter; setting it to zero retains all
    contaminants.

  \item \textbf{Raw-versus-finished reconciliation.} Where a system reports both
    raw (untreated) and finished (treated) samples in the same year and the
    finished-water mean is below the raw-water mean, the annual average is
    computed from finished samples only. This information is available in SYR3
    but not SYR2.

\end{enumerate}

After cleaning, sample-level records are collapsed to the
system\,$\times$\,contaminant\,$\times$\,year level. The annual concentration is
the mean of that year's (cleaned) samples; the share above the MCL is the share
of samples exceeding the applicable standard.

\subsection{Data-driven method detection limits}
\label{subsec:mdl}

Ravalli et al.\ impute the fallback LOD using EPA's published,
\emph{method-specific} maximum detection limits---a fixed externally specified
value per analytical method. We depart from this in one respect: rather than
hard-coding the EPA method MDLs, we derive each contaminant's fallback MDL
\emph{from the data}. Specifically, for every contaminant we take the median of
the reliable record-specific LODs actually reported in the Six-Year Review
(those in mg/L or $\mu$g/L and at most 5~$\mu$g/L); for contaminants whose
reported LODs all exceed that reliable threshold (e.g.\ nitrate), we instead use
the median of all reported mass-concentration LODs. This data-driven fallback is
used only for non-detect records that lack a usable record-specific LOD---most
importantly the SYR2 records, which do not carry a detection-limit field at all.

The data-driven approach has two advantages in our setting: it is reproducible
directly from the monitoring records without reconciling the multiple EPA
analytical methods in force across the 1998--2011 window, and it reflects the
detection limits laboratories actually achieved rather than the regulatory
ceiling. Its limitation is that the fallback MDL is estimated rather than
authoritative; the implementation therefore exposes an override table so that
official EPA method MDLs can be substituted directly if an exact replication of
Ravalli et al.\ is desired. In practice the imputed non-detect values
($\mathrm{MDL}/\sqrt 2$) are small relative to the MCL and to typical detected
concentrations, so the choice between data-driven and published MDLs has a
negligible effect on the estimated means.% (uses \mathrm below to avoid an amsmath dependency)

\subsection{Effect on the estimation sample}

Table~\ref{tab:6yr_cleaning_compare} summarizes the effect of the cleaning on the
downstream-CWS inorganic-chemical sample. Retaining imputed non-detects roughly
doubles the number of system\,$\times$\,year observations for each metal, while
the inclusion of these low (below-detection) values lowers the sample means.
Cadmium and mercury drop out under the detection-rate screen.

\begin{table}[htbp]
  \centering
  \caption{Inorganic-chemical concentrations before and after Ravalli-style
    cleaning (downstream-CWS regression sample, 1998--2011)}
  \label{tab:6yr_cleaning_compare}
  \begin{tabular}{lcccc}
    \toprule
    & \multicolumn{2}{c}{Before cleaning} & \multicolumn{2}{c}{After cleaning} \\
    \cmidrule(lr){2-3}\cmidrule(lr){4-5}
    Contaminant (mg/L) & Mean & $N$ & Mean & $N$ \\
    \midrule
    Arsenic   & 0.0041 & 451   & 0.0016 & 1{,}079 \\
    Nitrate   & 1.093  & 1{,}170 & 0.8367 & 1{,}486 \\
    Barium    & 0.0775 & 638   & 0.0305 & 799   \\
    Chromium  & 0.0078 & 410   & 0.0017 & 806   \\
    Selenium  & 0.0064 & 396   & 0.0030 & 796   \\
    Cadmium   & 0.0016 & 378   & \multicolumn{2}{c}{dropped ($<10\%$ detected)} \\
    Mercury   & 0.0003 & 375   & \multicolumn{2}{c}{dropped ($<10\%$ detected)} \\
    \bottomrule
  \end{tabular}
\end{table}
.
% ============================================================

\section{Cleaning of EPA Six-Year Review concentration data}
\label{sec:6yr_cleaning}

Measured contaminant concentrations are drawn from the U.S.\ EPA Six-Year
Review compliance-monitoring database, pooling the Second Six-Year Review
(SYR2, sample years 1998--2005) and Third Six-Year Review (SYR3, 2006--2011).
We clean these records following the procedure of Ravalli et al.\
(2022),\footnote{Ravalli F, Yu Y, Bostick BC, et al.\ Sociodemographic
inequalities in uranium and other metals in community water systems across the
USA, 2006--11: a cross-sectional study. \emph{Lancet Planet Health} 2022;
\textbf{6}: e320--30 (Supplementary appendix~2).} with one deliberate
departure described in Section~\ref{subsec:mdl}. All cleaning is performed on
the \emph{full national} sample of community water systems---before the analysis
sample is restricted to systems located downstream of a coal mine---so that
detection-rate screening, limit-of-detection imputation, and outlier removal
exploit the entire monitoring record rather than the small downstream subsample.

\subsection{Cleaning steps}

\begin{enumerate}

  \item \textbf{Community water systems only.} We retain records from community
    water systems (CWS) and drop transient and non-transient non-community
    systems, consistent with our focus on systems serving residential
    populations year-round. This removes roughly 23\% of raw monitoring records.

  \item \textbf{Unit normalization.} All mass-concentration readings are
    converted to mg/L and all radionuclide readings to pCi/L, so that records
    reported in $\mu$g/L, ng/L, or other units are placed on a common scale
    before averaging.

  \item \textbf{Limit-of-detection substitution.} Non-detect samples---which the
    database reports with a missing concentration---are assigned a value equal
    to the reported limit of detection (LOD) divided by $\sqrt{2}$, the standard
    convention used by the CDC and other federal agencies when summarizing
    environmental concentration data. A record's own reported LOD is used when it
    is reported in mg/L or $\mu$g/L and does not exceed 5~$\mu$g/L (the reliable
    range); otherwise the LOD is taken from a per-chemical method detection limit
    (MDL), as described in Section~\ref{subsec:mdl}. Imputing non-detects, rather
    than discarding them, is the main reason the analysis sample roughly doubles
    in size relative to a detections-only sample.

  \item \textbf{Gross-outlier removal.} Records reporting a concentration greater
    than 100 times the contaminant's maximum contaminant level (MCL) are dropped
    as likely data-entry errors.

  \item \textbf{Detection-rate screen.} For each contaminant we compute the
    national share of records above the LOD and retain only contaminants
    detected in more than 10\% of records. Among the inorganic chemicals this
    keeps arsenic, nitrate, barium, chromium, and selenium, and excludes cadmium
    (1.8\% detected), mercury (1.6\%), and thallium, whose concentrations are too
    rarely above detection to support a credible mean-concentration regression.
    The threshold is a tunable parameter; setting it to zero retains all
    contaminants.

  \item \textbf{Raw-versus-finished reconciliation.} Where a system reports both
    raw (untreated) and finished (treated) samples in the same year and the
    finished-water mean is below the raw-water mean, the annual average is
    computed from finished samples only. This information is available in SYR3
    but not SYR2.

\end{enumerate}

After cleaning, sample-level records are collapsed to the
system\,$\times$\,contaminant\,$\times$\,year level. The annual concentration is
the mean of that year's (cleaned) samples; the share above the MCL is the share
of samples exceeding the applicable standard.

\subsection{Data-driven method detection limits}
\label{subsec:mdl}

Ravalli et al.\ impute the fallback LOD using EPA's published,
\emph{method-specific} maximum detection limits---a fixed externally specified
value per analytical method. We depart from this in one respect: rather than
hard-coding the EPA method MDLs, we derive each contaminant's fallback MDL
\emph{from the data}. Specifically, for every contaminant we take the median of
the reliable record-specific LODs actually reported in the Six-Year Review
(those in mg/L or $\mu$g/L and at most 5~$\mu$g/L); for contaminants whose
reported LODs all exceed that reliable threshold (e.g.\ nitrate), we instead use
the median of all reported mass-concentration LODs. This data-driven fallback is
used only for non-detect records that lack a usable record-specific LOD---most
importantly the SYR2 records, which do not carry a detection-limit field at all.

The data-driven approach has two advantages in our setting: it is reproducible
directly from the monitoring records without reconciling the multiple EPA
analytical methods in force across the 1998--2011 window, and it reflects the
detection limits laboratories actually achieved rather than the regulatory
ceiling. Its limitation is that the fallback MDL is estimated rather than
authoritative; the implementation therefore exposes an override table so that
official EPA method MDLs can be substituted directly if an exact replication of
Ravalli et al.\ is desired. In practice the imputed non-detect values
($\mathrm{MDL}/\sqrt 2$) are small relative to the MCL and to typical detected
concentrations, so the choice between data-driven and published MDLs has a
negligible effect on the estimated means.% (uses \mathrm below to avoid an amsmath dependency)

\subsection{Effect on the estimation sample}

Table~\ref{tab:6yr_cleaning_compare} summarizes the effect of the cleaning on the
downstream-CWS inorganic-chemical sample. Retaining imputed non-detects roughly
doubles the number of system\,$\times$\,year observations for each metal, while
the inclusion of these low (below-detection) values lowers the sample means.
Cadmium and mercury drop out under the detection-rate screen.

\begin{table}[htbp]
  \centering
  \caption{Inorganic-chemical concentrations before and after Ravalli-style
    cleaning (downstream-CWS regression sample, 1998--2011)}
  \label{tab:6yr_cleaning_compare}
  \begin{tabular}{lcccc}
    \toprule
    & \multicolumn{2}{c}{Before cleaning} & \multicolumn{2}{c}{After cleaning} \\
    \cmidrule(lr){2-3}\cmidrule(lr){4-5}
    Contaminant (mg/L) & Mean & $N$ & Mean & $N$ \\
    \midrule
    Arsenic   & 0.0041 & 451   & 0.0016 & 1{,}079 \\
    Nitrate   & 1.093  & 1{,}170 & 0.8367 & 1{,}486 \\
    Barium    & 0.0775 & 638   & 0.0305 & 799   \\
    Chromium  & 0.0078 & 410   & 0.0017 & 806   \\
    Selenium  & 0.0064 & 396   & 0.0030 & 796   \\
    Cadmium   & 0.0016 & 378   & \multicolumn{2}{c}{dropped ($<10\%$ detected)} \\
    Mercury   & 0.0003 & 375   & \multicolumn{2}{c}{dropped ($<10\%$ detected)} \\
    \bottomrule
  \end{tabular}
\end{table}
.
% ============================================================

\section{Cleaning of EPA Six-Year Review concentration data}
\label{sec:6yr_cleaning}

Measured contaminant concentrations are drawn from the U.S.\ EPA Six-Year
Review compliance-monitoring database, pooling the Second Six-Year Review
(SYR2, sample years 1998--2005) and Third Six-Year Review (SYR3, 2006--2011).
We clean these records following the procedure of Ravalli et al.\
(2022),\footnote{Ravalli F, Yu Y, Bostick BC, et al.\ Sociodemographic
inequalities in uranium and other metals in community water systems across the
USA, 2006--11: a cross-sectional study. \emph{Lancet Planet Health} 2022;
\textbf{6}: e320--30 (Supplementary appendix~2).} with one deliberate
departure described in Section~\ref{subsec:mdl}. All cleaning is performed on
the \emph{full national} sample of community water systems---before the analysis
sample is restricted to systems located downstream of a coal mine---so that
detection-rate screening, limit-of-detection imputation, and outlier removal
exploit the entire monitoring record rather than the small downstream subsample.

\subsection{Cleaning steps}

\begin{enumerate}

  \item \textbf{Community water systems only.} We retain records from community
    water systems (CWS) and drop transient and non-transient non-community
    systems, consistent with our focus on systems serving residential
    populations year-round. This removes roughly 23\% of raw monitoring records.

  \item \textbf{Unit normalization.} All mass-concentration readings are
    converted to mg/L and all radionuclide readings to pCi/L, so that records
    reported in $\mu$g/L, ng/L, or other units are placed on a common scale
    before averaging.

  \item \textbf{Limit-of-detection substitution.} Non-detect samples---which the
    database reports with a missing concentration---are assigned a value equal
    to the reported limit of detection (LOD) divided by $\sqrt{2}$, the standard
    convention used by the CDC and other federal agencies when summarizing
    environmental concentration data. A record's own reported LOD is used when it
    is reported in mg/L or $\mu$g/L and does not exceed 5~$\mu$g/L (the reliable
    range); otherwise the LOD is taken from a per-chemical method detection limit
    (MDL), as described in Section~\ref{subsec:mdl}. Imputing non-detects, rather
    than discarding them, is the main reason the analysis sample roughly doubles
    in size relative to a detections-only sample.

  \item \textbf{Gross-outlier removal.} Records reporting a concentration greater
    than 100 times the contaminant's maximum contaminant level (MCL) are dropped
    as likely data-entry errors.

  \item \textbf{Detection-rate screen.} For each contaminant we compute the
    national share of records above the LOD and retain only contaminants
    detected in more than 10\% of records. Among the inorganic chemicals this
    keeps arsenic, nitrate, barium, chromium, and selenium, and excludes cadmium
    (1.8\% detected), mercury (1.6\%), and thallium, whose concentrations are too
    rarely above detection to support a credible mean-concentration regression.
    The threshold is a tunable parameter; setting it to zero retains all
    contaminants.

  \item \textbf{Raw-versus-finished reconciliation.} Where a system reports both
    raw (untreated) and finished (treated) samples in the same year and the
    finished-water mean is below the raw-water mean, the annual average is
    computed from finished samples only. This information is available in SYR3
    but not SYR2.

\end{enumerate}

After cleaning, sample-level records are collapsed to the
system\,$\times$\,contaminant\,$\times$\,year level. The annual concentration is
the mean of that year's (cleaned) samples; the share above the MCL is the share
of samples exceeding the applicable standard.

\subsection{Data-driven method detection limits}
\label{subsec:mdl}

Ravalli et al.\ impute the fallback LOD using EPA's published,
\emph{method-specific} maximum detection limits---a fixed externally specified
value per analytical method. We depart from this in one respect: rather than
hard-coding the EPA method MDLs, we derive each contaminant's fallback MDL
\emph{from the data}. Specifically, for every contaminant we take the median of
the reliable record-specific LODs actually reported in the Six-Year Review
(those in mg/L or $\mu$g/L and at most 5~$\mu$g/L); for contaminants whose
reported LODs all exceed that reliable threshold (e.g.\ nitrate), we instead use
the median of all reported mass-concentration LODs. This data-driven fallback is
used only for non-detect records that lack a usable record-specific LOD---most
importantly the SYR2 records, which do not carry a detection-limit field at all.

The data-driven approach has two advantages in our setting: it is reproducible
directly from the monitoring records without reconciling the multiple EPA
analytical methods in force across the 1998--2011 window, and it reflects the
detection limits laboratories actually achieved rather than the regulatory
ceiling. Its limitation is that the fallback MDL is estimated rather than
authoritative; the implementation therefore exposes an override table so that
official EPA method MDLs can be substituted directly if an exact replication of
Ravalli et al.\ is desired. In practice the imputed non-detect values
($\mathrm{MDL}/\sqrt 2$) are small relative to the MCL and to typical detected
concentrations, so the choice between data-driven and published MDLs has a
negligible effect on the estimated means.% (uses \mathrm below to avoid an amsmath dependency)

\subsection{Effect on the estimation sample}

Table~\ref{tab:6yr_cleaning_compare} summarizes the effect of the cleaning on the
downstream-CWS inorganic-chemical sample. Retaining imputed non-detects roughly
doubles the number of system\,$\times$\,year observations for each metal, while
the inclusion of these low (below-detection) values lowers the sample means.
Cadmium and mercury drop out under the detection-rate screen.

\begin{table}[htbp]
  \centering
  \caption{Inorganic-chemical concentrations before and after Ravalli-style
    cleaning (downstream-CWS regression sample, 1998--2011)}
  \label{tab:6yr_cleaning_compare}
  \begin{tabular}{lcccc}
    \toprule
    & \multicolumn{2}{c}{Before cleaning} & \multicolumn{2}{c}{After cleaning} \\
    \cmidrule(lr){2-3}\cmidrule(lr){4-5}
    Contaminant (mg/L) & Mean & $N$ & Mean & $N$ \\
    \midrule
    Arsenic   & 0.0041 & 451   & 0.0016 & 1{,}079 \\
    Nitrate   & 1.093  & 1{,}170 & 0.8367 & 1{,}486 \\
    Barium    & 0.0775 & 638   & 0.0305 & 799   \\
    Chromium  & 0.0078 & 410   & 0.0017 & 806   \\
    Selenium  & 0.0064 & 396   & 0.0030 & 796   \\
    Cadmium   & 0.0016 & 378   & \multicolumn{2}{c}{dropped ($<10\%$ detected)} \\
    Mercury   & 0.0003 & 375   & \multicolumn{2}{c}{dropped ($<10\%$ detected)} \\
    \bottomrule
  \end{tabular}
\end{table}
.
% ============================================================

\section{Cleaning of EPA Six-Year Review concentration data}
\label{sec:6yr_cleaning}

Measured contaminant concentrations are drawn from the U.S.\ EPA Six-Year
Review compliance-monitoring database, pooling the Second Six-Year Review
(SYR2, sample years 1998--2005) and Third Six-Year Review (SYR3, 2006--2011).
We clean these records following the procedure of Ravalli et al.\
(2022),\footnote{Ravalli F, Yu Y, Bostick BC, et al.\ Sociodemographic
inequalities in uranium and other metals in community water systems across the
USA, 2006--11: a cross-sectional study. \emph{Lancet Planet Health} 2022;
\textbf{6}: e320--30 (Supplementary appendix~2).} with one deliberate
departure described in Section~\ref{subsec:mdl}. All cleaning is performed on
the \emph{full national} sample of community water systems---before the analysis
sample is restricted to systems located downstream of a coal mine---so that
detection-rate screening, limit-of-detection imputation, and outlier removal
exploit the entire monitoring record rather than the small downstream subsample.

\subsection{Cleaning steps}

\begin{enumerate}

  \item \textbf{Community water systems only.} We retain records from community
    water systems (CWS) and drop transient and non-transient non-community
    systems, consistent with our focus on systems serving residential
    populations year-round. This removes roughly 23\% of raw monitoring records.

  \item \textbf{Unit normalization.} All mass-concentration readings are
    converted to mg/L and all radionuclide readings to pCi/L, so that records
    reported in $\mu$g/L, ng/L, or other units are placed on a common scale
    before averaging.

  \item \textbf{Limit-of-detection substitution.} Non-detect samples---which the
    database reports with a missing concentration---are assigned a value equal
    to the reported limit of detection (LOD) divided by $\sqrt{2}$, the standard
    convention used by the CDC and other federal agencies when summarizing
    environmental concentration data. A record's own reported LOD is used when it
    is reported in mg/L or $\mu$g/L and does not exceed 5~$\mu$g/L (the reliable
    range); otherwise the LOD is taken from a per-chemical method detection limit
    (MDL), as described in Section~\ref{subsec:mdl}. Imputing non-detects, rather
    than discarding them, is the main reason the analysis sample roughly doubles
    in size relative to a detections-only sample.

  \item \textbf{Gross-outlier removal.} Records reporting a concentration greater
    than 100 times the contaminant's maximum contaminant level (MCL) are dropped
    as likely data-entry errors.

  \item \textbf{Detection-rate screen.} For each contaminant we compute the
    national share of records above the LOD and retain only contaminants
    detected in more than 10\% of records. Among the inorganic chemicals this
    keeps arsenic, nitrate, barium, chromium, and selenium, and excludes cadmium
    (1.8\% detected), mercury (1.6\%), and thallium, whose concentrations are too
    rarely above detection to support a credible mean-concentration regression.
    The threshold is a tunable parameter; setting it to zero retains all
    contaminants.

  \item \textbf{Raw-versus-finished reconciliation.} Where a system reports both
    raw (untreated) and finished (treated) samples in the same year and the
    finished-water mean is below the raw-water mean, the annual average is
    computed from finished samples only. This information is available in SYR3
    but not SYR2.

\end{enumerate}

After cleaning, sample-level records are collapsed to the
system\,$\times$\,contaminant\,$\times$\,year level. The annual concentration is
the mean of that year's (cleaned) samples; the share above the MCL is the share
of samples exceeding the applicable standard.

\subsection{Data-driven method detection limits}
\label{subsec:mdl}

Ravalli et al.\ impute the fallback LOD using EPA's published,
\emph{method-specific} maximum detection limits---a fixed externally specified
value per analytical method. We depart from this in one respect: rather than
hard-coding the EPA method MDLs, we derive each contaminant's fallback MDL
\emph{from the data}. Specifically, for every contaminant we take the median of
the reliable record-specific LODs actually reported in the Six-Year Review
(those in mg/L or $\mu$g/L and at most 5~$\mu$g/L); for contaminants whose
reported LODs all exceed that reliable threshold (e.g.\ nitrate), we instead use
the median of all reported mass-concentration LODs. This data-driven fallback is
used only for non-detect records that lack a usable record-specific LOD---most
importantly the SYR2 records, which do not carry a detection-limit field at all.

The data-driven approach has two advantages in our setting: it is reproducible
directly from the monitoring records without reconciling the multiple EPA
analytical methods in force across the 1998--2011 window, and it reflects the
detection limits laboratories actually achieved rather than the regulatory
ceiling. Its limitation is that the fallback MDL is estimated rather than
authoritative; the implementation therefore exposes an override table so that
official EPA method MDLs can be substituted directly if an exact replication of
Ravalli et al.\ is desired. In practice the imputed non-detect values
($\mathrm{MDL}/\sqrt 2$) are small relative to the MCL and to typical detected
concentrations, so the choice between data-driven and published MDLs has a
negligible effect on the estimated means.% (uses \mathrm below to avoid an amsmath dependency)

\subsection{Effect on the estimation sample}

Table~\ref{tab:6yr_cleaning_compare} summarizes the effect of the cleaning on the
downstream-CWS inorganic-chemical sample. Retaining imputed non-detects roughly
doubles the number of system\,$\times$\,year observations for each metal, while
the inclusion of these low (below-detection) values lowers the sample means.
Cadmium and mercury drop out under the detection-rate screen.

\begin{table}[htbp]
  \centering
  \caption{Inorganic-chemical concentrations before and after Ravalli-style
    cleaning (downstream-CWS regression sample, 1998--2011)}
  \label{tab:6yr_cleaning_compare}
  \begin{tabular}{lcccc}
    \toprule
    & \multicolumn{2}{c}{Before cleaning} & \multicolumn{2}{c}{After cleaning} \\
    \cmidrule(lr){2-3}\cmidrule(lr){4-5}
    Contaminant (mg/L) & Mean & $N$ & Mean & $N$ \\
    \midrule
    Arsenic   & 0.0041 & 451   & 0.0016 & 1{,}079 \\
    Nitrate   & 1.093  & 1{,}170 & 0.8367 & 1{,}486 \\
    Barium    & 0.0775 & 638   & 0.0305 & 799   \\
    Chromium  & 0.0078 & 410   & 0.0017 & 806   \\
    Selenium  & 0.0064 & 396   & 0.0030 & 796   \\
    Cadmium   & 0.0016 & 378   & \multicolumn{2}{c}{dropped ($<10\%$ detected)} \\
    Mercury   & 0.0003 & 375   & \multicolumn{2}{c}{dropped ($<10\%$ detected)} \\
    \bottomrule
  \end{tabular}
\end{table}
